\originalchapter{研究生作业二：对偶与磁偶极}
题面据8.311作业2, 解答为AI独立推导、参照官方解答核对. 第7题为原卷选做题, 官方解答未给该题解.
\originalsection{角动量的局域守恒}
\begin{exercise}
原题1. 从动量守恒推出角动量密度$\vct l=\rr\times\vct g$及其通量, 核对局域角动量守恒式.
\end{exercise}
\begin{solution}
沿用第8章向外通量$-T_{ij}$: $\partial_tg_i-\partial_jT_{ij}=-f_i$. 乘$\epsilon_{abi}x_b$, 用乘积法则; 额外项$\epsilon_{aji}T_{ij}=0$因应力对称而消失. 定义$K_{aj}=-\epsilon_{abi}x_bT_{ij}$, 得$\partial_tl_a+\partial_jK_{aj}=-(\rr\times\vct f)_a$. 这就是原题对第二指标作叉积的分量含义. 整体角动量的改变来自边界通量及给物质的力矩; 非对称应力还需内部自旋角动量, 此处真空Maxwell应力对称.
\end{solution}
\originalsection{含磁荷的电磁对偶}
\begin{exercise}
原题2. (a) 将含磁荷Maxwell方程写成复场方程, 证明常相位旋转保持形式, 解释$\phi=\pi/2$; (b) 若每个粒子磁荷与电荷之比相同, 能否选择描述使磁荷消失?
\end{exercise}
\begin{solution}
SI采用$\nabla\cdot\bb=\mu_0\rho_m$、$\nabla\times\ee=-\partial_t\bb-\mu_0\jj_m$. 为让两源同量纲, 设$\vct F=\ee+ic\bb$, $\mathcal R=\rho_e+i\rho_m/c$, $\mathcal J=\jj_e+i\jj_m/c$. 逐式相加得到
$\nabla\cdot\vct F=\mathcal R/\eps$及$\nabla\times\vct F=(i/c)\partial_t\vct F+i\mathcal J/(\eps c)$. 同时乘$e^{-i\phi}$保留方程. 实部和虚部给$\ee'=\ee\cos\phi+c\bb\sin\phi$、$c\bb'=c\bb\cos\phi-\ee\sin\phi$, 源作同样二维转动. $\phi=\pi/2$时$\ee'=c\bb$、$c\bb'=-\ee$. 这是场及源的对偶旋转, 不是空间转动. 若普遍比值$\lambda=q_m/(cq_e)$相同, 选$\tan\phi=\lambda$便使所有磁荷为零, 电荷换成$q_e\sqrt{1+\lambda^2}$. 比值不一致时一个角度不能同时消去各磁荷.
\end{solution}
\originalsection{复场的双线性不变量}
\begin{exercise}
原题3. 解释$\vct F\cdot\vct F^*$、$\vct F\times\vct F^*$及对称二阶乘积, 并考察对偶变换.
\end{exercise}
\begin{solution}
由上一节定义, 能量$u=\eps\vct F\cdot\vct F^*/2$, 动量密度$\vct g=i\eps\vct F\times\vct F^*/(2c)=\eps\ee\times\bb$. 对称张量$\eps(F_iF_j^*+F_i^*F_j)/2=\eps E_iE_j+B_iB_j/\mu_0=T_{ij}+u\delta_{ij}$. 因$\vct F$及共轭的相位相反, 各乘积均不变, 因而能量、动量和应力都不因对偶改变. 原高斯系数必须连同$F$定义一起换算.
\end{solution}
\originalsection{任意小电流环的力和力矩}
\begin{exercise}
原题4. (a) 任意形状闭合电流环在均匀磁场中求总力和力矩; (b) 外场在环尺度上弱变化, 求首阶合力并用磁矩表示.
\end{exercise}
\begin{solution}
(a) $\vct F=I\oint\dd\vct l\times\bb=0$. 记$m_i=I\epsilon_{ijk}\oint x_j\dd x_k/2$, 将$\rr\times(\dd\rr\times\bb)$展开并用$\oint\dd(x_ix_j)=0$, 得$\vct\tau=\mm\times\bb$. (b) 展开$B_k(\rr)=B_k(0)+x_l\partial_lB_k$, 再用闭环一阶矩反对称性, 得$F_i=\partial_i(\mm\cdot\bb)$, 其中$\mm$固定, $\nabla\cdot\bb=0$. 因而$\vct F=\nabla(\mm\cdot\bb)$; 若写$(\mm\cdot\nabla)\bb$, 还须在环所在区$\nabla\times\bb=0$, 否则两式有$\mm\times(\nabla\times\bb)$差别.
\end{solution}
\originalsection{偶极场与相互作用能}
\begin{exercise}
原题5. (a) 从小电流环远场推出磁偶极场; (b) 求两固定磁偶极的相互作用能.
\end{exercise}
\begin{solution}
第6章电流矩展开已给$\vct A=\mu_0\mm\times\rr/(4\pi r^3)$. 在源外取旋度, 用$\partial_i(x_j/r^3)=\delta_{ij}/r^3-3x_ix_j/r^5$, 得$\bb=\mu_0[3\nn(\mm\cdot\nn)-\mm]/(4\pi r^3)$. 对磁矩$\mm_2$, 转动做功给$U=-\mm_2\cdot\bb_1$, 因而$U=\mu_0[\mm_1\cdot\mm_2-3(\mm_1\cdot\nn)(\mm_2\cdot\nn)]/(4\pi R^3)$. 这里$R=|\rr_2-\rr_1|$远大于各环尺寸, 未含偶极自身能.
\end{solution}
\originalsection{旋转球壳与场角动量}
\begin{exercise}
原题6. 半径$a$、总电荷$q$的均匀带电球壳以$\omega\hat{\vct z}$旋转. (a) 求电流与磁矩; (b,c) 联立边界条件求内外磁场并与磁矩核对; (d) 求场角动量.
\end{exercise}
\begin{solution}
(a) 面电流$\vct K=q\omega\sin\theta\hat{\vct\phi}/(4\pi a)$, 体密度$\jj=\vct K\delta(r-a)$, 由$\mm=\int\rr\times\jj\dd^3x/2$得$\mm=q\omega a^2\hat{\vct z}/3$.

(b,c) 第6章已实际解出磁化球边界方程: 取$M=q\omega/(4\pi a)$, 内场$B_0=2\mu_0M/3=\mu_0q\omega/(6\pi a)$, 外场$\mu_0[3\nn(\mm\cdot\nn)-\mm]/(4\pi r^3)$. $B_n$连续、$\nn\times(\bb_{\rm out}-\bb_{\rm in})=\mu_0\vct K$逐项满足.

(d) 球内电场零, 外部$E_r=q/(4\pi\eps r^2)$与$B_\theta=\mu_0m\sin\theta/(4\pi r^3)$给$g_\phi=\mu_0qm\sin\theta/(16\pi^2r^5)$. 故$L_z=\int r\sin\theta g_\phi\dd^3x=\mu_0qm/(6\pi a)=\mu_0q^2\omega a/(18\pi)$. 此角动量与本科转动球壳的附加惯量一致, 不是实心球系数.
\end{solution}
\originalsection{电荷与磁单极的角动量（选做）}
\begin{exercise}
原题7（选做）. 电荷$q$在$\vct R/2$, 磁荷$g$在$-\vct R/2$. (a) 写动量密度并表示成旋度; (b) 求总场角动量, 证明只依赖$\vct R$方向, 解释半经典Dirac量子化.
\end{exercise}
\begin{solution}
记$u_e=1/|\rr-\vct R/2|$, $u_m=1/|\rr+\vct R/2|$, $C=\mu_0qg/(16\pi^2)$. 则$\vct g=C\nabla u_e\times\nabla u_m=C\nabla\times(u_e\nabla u_m)$, 无散（奇点按挖洞极限）. 分部积分$\int\rr\times\nabla\times\vct A=2\int\vct A$给$\vct L=2C\int u_e\nabla u_m$; 无穷远的领先径向项不贡献叉积边界项.

为计算它, 考察正则化卷积$J(R)=\int u_eu_m\dd^3x$并减去与$R$无关的发散常量. $\nabla_R^2J=-4\pi/R$、球对称使$(R^2J')'=-4\pi R$, 积分得$J'=-2\pi+C_1/R^2$. $R\to0$无$1/R$奇项令$C_1=0$, 即$J'(R)=-2\pi$. 对磁荷位置微分, $\int u_e\nabla u_m=-2\pi\hat{\vct R}$. 因而$\vct L=-\mu_0qg\hat{\vct R}/(4\pi)$, 与距离无关. 符号取决于$\vct R$由磁荷指向电荷的约定. 以角动量投影的半整数条件作半经典解释, $\mu_0qg/(4\pi)=n\hbar/2$, 对应原高斯制$eg=n\hbar c/2$. 这不是磁荷存在的实验结论, 也不是完整量子场论证明.
\end{solution}
